\section{The Exact Response Kernel and Best Possible Scalar Bounds}\label{app:kernel}

The response kernel measures how strongly the averaged update reacts to a
small parameter change. Fix one coordinate, omit its index, and let $R_j=|e_{-j}|$,
$V=\sum_jd_jR_j^2$, $D_V=\sqrt V$. Since independent time samples have
$\Pp(R>0)>0$, both $V$ and each $V_{-j}=V-d_jR_j^2$ are positive
almost surely. Differentiate the frozen quotient in direction $k$:
\[
 \partial_kf_i(0,Z)=
 \frac{\sum_jw_jA_{ik,-j}}{D_V+\epsilon}
 -\frac{(\sum_lw_le_{i,-l})(\sum_jd_je_{i,-j}A_{ik,-j})}
 {D_V(D_V+\epsilon)^2}.
\]
For $l\ne j$, condition on all other times and $|e_{i,-l}|$. The
corresponding product contains $e_{i,-l}$ but not the curvature from that
same sample. Its conditional expectation is zero because the noise
marginal is symmetric. No sign symmetry conditional on that absent
curvature is needed. The surviving same-sample terms are
\[
 \sum_j w_j A_{ik,-j}
 \frac{V-d_jR_j^2+\epsilon\sqrt V}{\sqrt V(\sqrt V+\epsilon)^2}.
\]
Condition each on its own $(A_{ik,-j},R_j)$ and average other times.
This gives \eqref{eq:rowkernel} and \eqref{eq:psi}.

To justify differentiating inside the expectation, use the weighted
square-sum norm on gradient histories $\|x\|_d^2=\sum_jd_jx_j^2$. The numerator
$n(x)=\sum_jw_jx_j$ has representing vector of norm $C_*$. The map
$T(x)=n(x)/(\|x\|_d+\epsilon)$ is continuously differentiable ($C^1$), including derivative
$n/\epsilon$ at zero, with a bounded Lipschitz derivative: its change is at
most a constant times the change in the history. Bounded affine gradient histories
therefore allow differentiation under expectation.
The own-time factor is at most $1/(D_V+\epsilon)\leq1/\epsilon$;
its sum is uniformly bounded by the summable weights.

For $v>0$, both
\[
 \frac{v}{x(x+\epsilon)^2},\qquad
 \frac{\epsilon}{(x+\epsilon)^2},\qquad x=\sqrt{v+d_jr^2},
\]
are strictly decreasing with $r\geq0$. The $j=0$ term has positive
weight, including when $\alpha=0$. Dominated convergence proves continuity;
strictness of this one term proves strict monotonicity of $\Psi$.
Because $D_V\leq R_{\max}$ after each substituted $r\leq R_{\max}$,
\[
 0<\frac{\epsilon}{(R_{\max}+\epsilon)^2}
       \leq\Psi(r)\leq\frac1\epsilon.
\]

\paragraph{Why scalar drift points toward zero.}
For a numerator summand at age $j$, condition on all other times and put
$J=\sum_{k\ne j}d_kg_{-k}(u)^2$. The function
$x/(\sqrt{d_jx^2+J}+\epsilon)$ is odd increasing. For $u>0$ and $H_j>0$,
its value at $H_ju+e_j$ is at least its value at $e_j$ pointwise, including
when $H_j$ depends on the sign. The latter averages to zero by marginal
symmetry and independence of other times. Reverse the comparison for $u<0$.
Summing weights yields $uF(u)\geq0$, hence
$uG(u)\geq\lambda u^2$. The field need not be odd when curvature can depend on the noise sign.
Lemma~\ref{lem:regular-transfer} below gives the stated
stationary law and moments after positivity of the origin variance.

\paragraph{Fluctuations at zero.}
At the origin the action depends only on the common noise law. Conditional
on its magnitudes, each nonzero sign is fair and independent across time.
For a scalar lag $l\geq0$,
\[
 \E[U_0U_l\mid |e|]=
 \frac{\alpha^lQ_0}{(\sqrt{V_0}+\epsilon)(\sqrt{V_l}+\epsilon)},
 \qquad Q_0=\sum_jw_j^2e_{-j}^2.
\]
This is nonnegative. It is bounded by
$[(1-\alpha)^2/(1-\beta)](\alpha/\rho)^l$, and its variance term is
strictly positive. Thus $0<D_\epsilon<\infty$. The martingale-corrector
variance agrees with this covariance sum by the bounded telescoping
decomposition and absolute summability.

\paragraph{Best and worst scalar dependence.}
For nonnegative bounded $H,W$, write their product as an integral of
joint tail probabilities:
\[
 \E HW=\int_0^\infty\!\int_0^\infty
                   \Pp(H>x,W>y)\,dx\,dy.
\]
Writing $A_x=\Pp(H>x)$ and $B_y=\Pp(W>y)$, intersection probabilities
lie between $\max(0,A_x+B_y-1)$ and $\min(A_x,B_y)$. Opposite/equal
rank pairings attain these bounds simultaneously: pair larger values
with smaller values for the lower bound, or larger with larger for the
upper bound. Apply this to
$W=\Psi(R)$, whose order reverses $R$. This gives the smallest and largest response slopes and, because risk
decreases with the slope, the risk endpoints.
Fair independent signs make both extremizers admissible in the signed
class. Mixtures preserve both marginals and $\E He=0$, fill the entire
susceptibility interval linearly, and fill the risk interval continuously.
If neither marginal is constant, the intersection bounds differ on a set
of positive area, so the two endpoints differ strictly.

Finally, $h_{\min}\E\Psi(R)\leq\chi\leq h_{\max}\E\Psi(R)$ yields
\[
 \frac{C_{\max}}{C_{\min}}
 =\frac{\lambda+\chi_{\max}}{\lambda+\chi_{\min}}
 \leq\frac{\lambda+h_{\max}\E\Psi(R)}
             {\lambda+h_{\min}\E\Psi(R)}
 \leq\frac{h_{\max}}{h_{\min}}.
\]
This bounds curvature spread across samples. It differs from the condition
number of one scalar quadratic, which is always one. The bound does not
say that every choice of marginals can attain this ceiling.
